
\subsubsection{Trustworthiness Benchmarks}

TruthfulQA measures factual accuracy across 817 questions designed to test whether models reproduce common misconceptions. BBQ (Bias Benchmark for QA) evaluates stereotype bias through ambiguous-context questions spanning nine demographic categories. For adversarial robustness, benchmarks include AdvGLUE and ANLI, which test model resilience to adversarially-constructed inputs. These benchmarks have enabled systematic evaluation of individual trustworthiness dimensions.

\subsubsection{Multi-Dimensional Evaluation Frameworks}

Recent frameworks evaluate models across multiple trustworthiness dimensions simultaneously. TrustLLM and DecodingTrust assess models on six principles including truthfulness, safety, fairness, robustness, privacy, and machine ethics. These frameworks document that models exhibit varying performance profiles across dimensions, with no model performing uniformly well. However, these frameworks focus on comprehensive measurement rather than characterizing how dimensions interact under interventions.

\subsubsection{Multi-Task Learning}

The multi-task learning literature demonstrates that tasks sharing neural representations can exhibit negative transfer, where improving one task degrades another. Gradient-based interference analysis shows that conflicting gradient directions in shared parameters create optimization trade-offs. This work differs by investigating post-training interventions rather than joint optimization during training, and by characterizing trustworthiness dimensions as emergent capabilities rather than explicit training targets.

\subsubsection{Gap in Prior Work}

No prior work systematically characterizes which trustworthiness dimension pairs exhibit trade-offs versus independence under targeted post-training interventions, or whether patterns replicate across model architectures. This study addresses that gap through perturbation-based correlation analysis.
